\section{Introduction}
\label{sec:intro}

Robotic garment manipulation is a benchmark for dexterous bimanual control. Unlike rigid object manipulation, handling textiles requires reasoning about infinite-dimensional configuration spaces and managing complex self-occlusions. The ICRA 2026 LeHome Challenge \cite{lehome2026} provides a rigorous platform for evaluating these capabilities, requiring robots to fold long-sleeved tops, short-sleeved tops, long pants, and shorts in a high-fidelity simulation environment \cite{mittal2023orbit}.

This paper documents our systematic exploration of various visuomotor policies and visual representations to solve the bimanual folding task. Rather than focusing on a single architecture, we treat this project as an empirical study of the most promising strategies in modern imitation learning. Our journey followed a deliberate progression: beginning with Diffusion Policy \cite{chi2023diffusionpolicy}, we then explored the emerging class of Vision-Language-Action (VLA) models, including SmolVLA \cite{smolvla2025}, $\pi_{0.5}$ \cite{pi05_2026}, and X-VLA \cite{xvla2026}. Ultimately, we transitioned to Action Chunking with Transformers (ACT) \cite{zhao2023act}. We found ACT to be superior for this task due to its inherent data efficiency, its specific design for bimanual coordination, and its ability to handle fine-grained, contact-rich manipulation. Furthermore, with only 80 million parameters, ACT is significantly more lightweight than the foundation-scale VLA models, making it ideal for rapid iteration.

We evaluate the role of geometric priors by providing colorized depth maps as an additional input, the impact of data augmentations on cross-style generalization, and the utility of Parameter-Efficient Fine-Tuning (PEFT) for adapting large-scale models. We also perform a sensitivity analysis on temporal parameters, specifically the action chunk size and the number of execution steps, and compare the performance of ACT when using a DINOv2 backbone versus the standard ResNet-18.

Our results provide a comprehensive view of the strengths and weaknesses of current architectures. While small VLA models demonstrate strong semantic understanding, large VLA models struggle when fine-tuned on the relatively small datasets typical of specialized manipulation tasks. We find that transformer-based temporal modeling in ACT, combined with geometric priors and robust augmentations, provides the most reliable performance for the specific constraints of the LeHome Challenge. We report detailed results across all garment categories, achieving a final success rate of 63.54\%.

In summary, our contributions include:
\begin{itemize}[leftmargin=*,noitemsep]
    \item A comparative study of multiple visuomotor architectures (DP, SmolVLA, $\pi_{0.5}$, X-VLA, ACT) for bimanual garment folding.
    \item An evaluation of mid-level visual representations (colorized depth) and their impact on reducing appearance-based overfitting.
    \item A systematic analysis of data augmentation, PEFT, and backbone selection (DINOv2) for improving task performance and generalization.
    \item A detailed breakdown of success rates per garment category and an analysis of critical temporal hyperparameters.
\end{itemize}
